%% sections/introduction.tex -- Phase 7.1. Register: CLAUDE.md S9.2 (no em/en dashes as prose
%% punctuation; one idea per sentence; numbers with meaning; paragraph openers carry the argument).
%% RQ1 abstract/intro fence: never "significant across all datasets"; enumerate cells instead.
%% QWK is deliberately NOT used in this section: its first use, with the mandatory calibration
%% sentence, is in Method (PHASES 7.1).

\section{Introduction}
\label{sec:intro}

Large language models (LLMs) are moving from demonstrations to deployed graders in computer science
(CS) courses. The promise for learning analytics is well documented: rubric-conditioned LLMs approach
human raters on well-specified short-answer tasks, and they return criterion-level feedback at a scale
no teaching team can match~\cite{sessler_canai_2025,henkel_roars_2025}. The risk is equally well
documented. Grading errors distort how students revise their work~\cite{li_am_2023}, repeated scoring
of the same answer is unstable~\cite{pack_large_2024,moazzez_assessing_2024}, and agreement with
teachers varies sharply between models and prompts~\cite{tate_can_2024}. An instructor who adopts an
LLM grader is therefore not making one decision but several: whether the model should reason at length
before scoring, what evaluation guidance to place in the context, whether to grade a whole exam in one
call or question by question, whether to score a rubric criterion by criterion or holistically, and
which model family to run or pay for.

The literature answers these questions only in fragments. Reasoning is the clearest example. Extended
thinking improves general LLM-as-a-judge accuracy~\cite{jayarao_explicit_2025}, but its effect is task-
and architecture-dependent~\cite{larionov_deepseek_2025}, long reasoning traces carry verbose,
redundant computation~\cite{sui_stop_2025}, and no peer-reviewed study isolates its effect on
rubric-bounded educational grading. Model comparisons exist for programming and course
deployments~\cite{jukiewicz_comparison_2026,policar_automated_2025}, but they rank systems under one
fixed prompt configuration, so they cannot say how the ranking moves when the configuration does.
Guidance, scope, decomposition, and conversation state have each been touched as point solutions
(Section~\ref{sec:related}), never crossed under one protocol. What is missing is the practical
synthesis: for a given CS assessment task, which configuration should you use, and what does it cost?

This paper studies the configuration space systematically and delivers the synthesis as a decision
guide. We cross six design dimensions: extended reasoning (on or off), evaluation guidance (none
versus a rubric or reference answer), grading scope (whole exam versus question by question), rubric
decomposition (holistic versus criterion by criterion), model family, and conversation state (a clean
conversation per answer versus shared history). Three open-weight model families (Qwen3.5,
DeepSeek-V4-Flash, GLM-5.1) expose a within-family reasoning toggle, so the reasoning effect is
isolated from model identity, a confound in prior comparisons; a closed frontier model (GPT-5.1)
serves as a reference point. The study runs on four datasets in two domains. Three are public and give cross-population external
validity: the Mohler~\cite{mohler_text_2009,mohler_learning_2011} and
SemEval-2013 Task~7~\cite{dzikovska_semeval_2013} corpora of graded computer-science short answers,
and the Rubric Is All You Need corpus of rubric-graded Java code~\cite{pathak_rubric_2025}. The
fourth, PT-CS, is an export of a real Portuguese university grading deployment built by the first
author; we use it to test whether the public-dataset guidance transfers to a deployment context, and
because its grades are of uneven reliability we curate a teacher-reviewed subset and set out the
accompanying safeguards in Section~\ref{sec:threats}.

Six research questions structure the study:
\begin{itemize}
  \item \textbf{RQ1.} How does extended reasoning affect grading quality, consistency, and cost, and
        does the effect differ between code and short answers?
  \item \textbf{RQ2.} How much does evaluation guidance (none versus a rubric or reference answer)
        change quality and consistency?
  \item \textbf{RQ3.} How do grading scope (whole exam versus question by question) and rubric
        decomposition (holistic versus criterion by criterion) affect agreement with the final grade?
  \item \textbf{RQ4.} How do current-generation models compare as graders, and what is the
        cost-quality trade-off between open-weight and closed frontier models?
  \item \textbf{RQ5.} Does conversation state (clean versus shared history) affect consistency and
        fairness?
  \item \textbf{RQ6.} Does the configuration guidance found on public datasets transfer to a real
        Portuguese deployment context?
\end{itemize}

This synthesis takes the form of a framework, not a leaderboard. Rather than crowning one best grader,
we build a decision guide that maps the characteristics of a grading task, its domain, its rubric, its
budget, and the reliability of the reference grades available, to a recommended configuration, the
cost that choice implies, and the conditions under which it applies. Section~\ref{sec:results} reports
the controlled comparisons behind the guide, and Section~\ref{sec:framework} presents it.

The contribution is fourfold. First, a configuration decision guide for LLM grading in CS, keyed by
observable task and model behaviour rather than by model identity. Second, the first
grading-specific, within-family isolation of the reasoning effect, measured jointly on agreement,
consistency, and cost. Third, a strategy for deployment grades of uneven reliability: an auditable
criterion that isolates a teacher-reviewed subset, so the deployment conclusions rest on a reference we
can trust. Fourth, a released harness, prompt templates, and configuration matrices that make the
protocol reproducible.

The paper is organised as follows. Section~\ref{sec:related} positions the study in the
learning-analytics and LLM-grading literature. Section~\ref{sec:method} describes the design space,
datasets, grading harness, and experimental protocol. Section~\ref{sec:results} reports the
experiments and their results by research question. Section~\ref{sec:discussion} discusses the
cross-cutting findings and Section~\ref{sec:framework} presents the configuration framework.
Section~\ref{sec:threats} states the threats to validity, and Section~\ref{sec:conclusion} draws the
study's conclusions and outlines future work.
